\documentclass[11pt]{article}
\usepackage[margin=2.6cm]{geometry}
\usepackage[T1]{fontenc}
\usepackage[utf8]{inputenc}
\usepackage{lmodern}
\usepackage{amsmath,amssymb}
\usepackage{array,longtable}
\usepackage{listings}
\usepackage{xcolor}
\usepackage[hidelinks]{hyperref}
\usepackage{parskip}

\lstset{
  basicstyle=\ttfamily\small,
  breaklines=true,
  columns=fullflexible,
  keepspaces=true,
  frame=single,
  framerule=0.2pt,
  xleftmargin=0.6em,
  aboveskip=0.7em,
  belowskip=0.7em,
  literate={—}{{---}}1 {°}{{\ensuremath{^\circ}}}1 {≈}{{\ensuremath{\approx}}}1
           {–}{{--}}1 {’}{{'}}1
}

\newcommand{\kw}[1]{\texttt{\uppercase{#1}}}
\newcommand{\code}[1]{\texttt{#1}}

\title{SolveIt\\[2pt]\large The Complete Solution Guide}
\author{rustSolveIt, on pure-Rust SUNDIALS 7.8.0}
\date{17 August 2026}

\begin{document}
\maketitle

\begin{center}
\emph{Written for a reader who has never used this program, has never run
a physics simulator, and does not write Rust. Every term is defined the
first time it is used. Every number printed below was produced by the
program on the machine that wrote this file; nothing is quoted from
memory.}
\end{center}

\tableofcontents
\newpage

\section{What this is, in one page}

\textbf{SolveIt} simulates classical mechanics and one-, two- and
three-dimensional quantum mechanics. You describe a situation --- some
bodies, some fields, maybe a rigid box to put them in --- and it tells
you what happens next, to as many digits as the arithmetic allows.

You talk to it by typing short commands, one per line:

\begin{lstlisting}
In[1]:= new sphere { mass = 2, radius = 0.5, position = [0, 10, 0],
                     velocity = [1, 0, 0] }
Out[1]= obj0
In[2]:= set system.gravity = [0, -9.81, 0]
In[3]:= step 1
Out[3]= t = 1 (advanced by 1, 12 solver steps)
In[4]:= get obj0.position
Out[4]= [1, 5.095000000000006, 0]
\end{lstlisting}

That is the whole interface. No configuration files, no build system to
learn, no graphics API.

\textbf{The one thing that makes this different from a game engine.}
Game physics is allowed to be approximately right --- it only has to
look plausible for the next sixteen milliseconds. This program's job is
to be \emph{measurably} right. Every trajectory it prints comes out of a
production-grade differential-equation solver with genuine error
control, and every claim in the documentation is a number you can
reproduce by running the command underneath it.

Concretely, and these are the actual measured values:

\begin{center}
\begin{tabular}{p{7.2cm}p{7.2cm}}
\textbf{claim} & \textbf{measured}\\\hline
the outer solar system, integrated for 1{,}369 years & Pluto's position
matches the reference to 8 decimal places; total energy drifts by 7.8
parts in $10^{7}$\\
a torque-free tumbling body & angular momentum changes by
\textbf{exactly zero}\\
a ball dropped on a plate & the moment of impact is right to 1 part in
$10^{16}$\\
three shapes rattling in a rigid box, 36 collisions & energy conserved
to 3 parts in $10^{16}$\\
Kepler's third law across four orbits & $T^{2}/a^{3}$ identical to the
last bit for all four\\
\end{tabular}
\end{center}

\subsection{The rules the code lives under}

\begin{enumerate}
\item \textbf{Everything is integrated by SUNDIALS.} There is no
      hand-written stepper anywhere in the program --- not in the
      examples, not in the graphics playback, not as a ``fast path''.
\item \textbf{No \code{unsafe} code.} Rust's memory-safety checks are on
      everywhere, with no escape hatches.
\item \textbf{No dependencies at all.} Not one package from the
      internet. The solver library, the web server behind the graphics
      window, the WebSocket protocol, the SHA-1 implementation, the
      special-function library --- all of it is in this repository.
\item \textbf{No compiler warnings.} The build is clean, and it is set
      up to \emph{fail} if it stops being clean.
\end{enumerate}

Rule 3 has a consequence worth spelling out: \code{git clone},
\code{cargo build}, done. No network access during the build, nothing to
go stale, nothing that can be pulled out from under you.

\section{Installing and running it}

You need Rust. Then:

\begin{lstlisting}
git clone https://github.com/once-ere/rustSolveIt_macos-silicon_SUNDIALS_7_8_0.git
cd rustSolveIt_macos-silicon_SUNDIALS_7_8_0
cargo run
\end{lstlisting}

The last command builds everything and drops you at the \code{In[1]:=}
prompt. Type \kw{help} for the reference card, \kw{quit} to leave.

\begin{center}
\begin{tabular}{p{3.4cm}p{7.4cm}p{3.6cm}}
\textbf{how} & \textbf{command} & \textbf{when}\\\hline
interactive notebook & \code{cargo run} & you are exploring\\
batch script & \code{cargo run -p posim -{}- -{}-script my\_run.posim} &
you want it reproducible\\
dynamic notebook & \code{cargo run -p posim -{}-release -{}- -{}-notebook
dynamic\_notebooks/kepler\_orbit.posim} & run a saved session, open its
graphics window, keep typing\\
machine mode & \code{cargo run -p posim -{}- -{}-machine} & another
program is driving\\
\end{tabular}
\end{center}

A ``script'' is just a text file of the same commands you would type.
Comments start with \code{\#}. Every example in \S\ref{sec:examples} is a
script in \code{scripts/solveit/}.

Checking your build: \code{cargo test -{}-workspace} --- 605 tests should
pass.

\section{The five things you need to know}

\subsection{\kw{run} takes a duration, not a destination}

\begin{lstlisting}
In[1]:= run 1.7 steps 170        # advances BY 1.7, so t is now 1.7
In[2]:= run 1.19 steps 120       # advances BY 1.19, so t is now 2.89
\end{lstlisting}

If you want to stop at an absolute time, subtract. This trips up
everyone once. \kw{step} \emph{dt} is the same thing with one output
point.

\subsection{Gravity is on by default, and it is between \emph{bodies}}
\label{sec:gravity}

\code{system.g\_constant} defaults to \textbf{1}, which means every pair
of bodies attracts every other pair. That is what you want for a solar
system and emphatically not what you want for three blocks in a box: two
touching surfaces are at nearly zero separation, and the attraction
between them is nearly singular.

If you are studying collisions, turn it off:
\code{set system.g\_constant = 0}.

The difference is not subtle. The same three-shapes-in-a-box scenario
conserves energy to \textbf{3 parts in $10^{16}$} with $G=0$, and drifts
by \textbf{3.2\,\%} with $G=1$. Neither number is a bug --- they are
different physical systems. \code{system.uniform\_gravity} is the
separate thing you usually mean by ``gravity'': a constant downward
field, like a laboratory.

\subsection{Objects are numbered, and deleting renumbers them}

\code{obj0}, \code{obj1}, \dots\ in creation order. \kw{del} 1 removes
\code{obj1} \textbf{and renumbers everything after it}. Give things
names instead: \code{new sphere as earth \{ \dots\ \}}, then
\code{get earth.position}. Names survive deletions of other objects.

\subsection{Collisions must be switched on}

\kw{collide on}. Without it, bodies pass through each other. This is a
deliberate default: collision detection changes how the solver picks its
steps, and a great many problems (orbits, fields, quantum) never need
it.

\subsection{$|dE/E|$ is printed for free, and you should read it}

Every \kw{run} reports the relative change in total energy. For a closed
system this should be tiny. If it is not, either the physics genuinely
does not conserve energy (a magnetic torque, an inelastic bounce, a
driven potential --- see examples 4 and 8) or something is wrong with
your setup. It is the cheapest sanity check you will ever get.

\section{What the program is made of}

\begin{center}
\begin{tabular}{p{3.2cm}p{11.4cm}}
\textbf{piece} & \textbf{what it does}\\\hline
\code{posim} & the command language: reads your line, checks it,
compiles it to a tiny stack program, runs it. Also the graphics window
and the machine-mode protocol.\\
\code{physical\_object} & the physics: one \code{physical\_object} type
that is simultaneously a point mass, a rigid body and a shaped,
collidable solid. Plus all the time integration.\\
\code{quantum} & 1-D, 2-D and 3-D quantum mechanics: bound states,
wavepacket propagation, tunnelling, absorbing boundaries.\\
\code{special\_functions} & Bessel, Legendre, Hankel, Airy, gamma,
Wigner symbols, quadrature, eigenproblems.\\
\code{sundials\_rs} & the differential-equation solvers (\S\ref{sec:engine}).\\
\code{jupyter} & a JupyterLab kernel.\\
\end{tabular}
\end{center}

\subsection{The one type that does everything}

The original 2002 SolveIt had three separate types: a point particle, a
rigid body, and a 3-D rigid body. Each had its own integrator, its own
bugs, and its own idea of what ``position'' meant. This version has
\textbf{one} type, \code{physical\_object}, which is the union of all
three. A point mass is just a \code{physical\_object} whose shape is
\kw{point} and whose inertia tensor is zero.

That sounds like a detail. It is why a point mass, a spinning cuboid and
a dumbbell can all be in the same collision and the same conservation
sum, without a single special case.

\subsection{The state vector}

Every object carries \textbf{13 numbers}:
\[
[\;\text{position (3)}\;|\;\text{momentum (3)}\;|\;
  \text{orientation quaternion (4)}\;|\;\text{angular momentum (3)}\;]
\]
Momentum rather than velocity, and angular \emph{momentum} rather than
angular velocity, both on purpose: those are the quantities that are
conserved, and the ones that behave correctly when you change a mass
mid-run. Quaternions are written \textbf{w first}: $[w,x,y,z]$.

\section{The engine: SUNDIALS 7.8.0 in pure Rust}\label{sec:engine}

\subsection{What SUNDIALS is}

\textbf{SUNDIALS} --- SUite of Nonlinear and DIfferential/ALgebraic
equation Solvers --- is Lawrence Livermore National Laboratory's library
of differential-equation solvers, in production use in physics and
engineering codes since the 1990s. It is written in C.

\code{sundials\_rs/} in this repository is a \textbf{line-for-line
translation of SUNDIALS 7.8.0 into pure Rust}. Not a wrapper, not
bindings: there is no C compiler involved, no shared library to link, no
FFI boundary. It carries the same control flow, the same constants, the
same tolerances, and --- a subtler point that matters a great deal ---
the same \emph{order of arithmetic operations}, because floating-point
addition is not associative and changing the order changes the answer in
the last digits.

\subsection{Which solver runs when}

\begin{center}
\begin{tabular}{p{4.2cm}p{10.4cm}}
\textbf{you type} & \textbf{what runs}\\\hline
\kw{method adams} (default) & \textbf{CVODE Adams--Moulton}: variable
order, variable step size, Newton iteration, dense finite-difference
Jacobian. The general-purpose choice.\\
\kw{method bdf} & \textbf{CVODE BDF}: backward differentiation formulas,
for \emph{stiff} problems --- ones with two very different timescales,
like a fast magnetic gyration inside a slow drift.\\
\kw{method sprk} \emph{table} [\emph{dt}] & \textbf{ARKODE SPRKStep}: a
\emph{symplectic} partitioned Runge--Kutta method at a fixed step. Only
valid for separable systems, and worth it when you have one.\\
\end{tabular}
\end{center}

``Adaptive'' means the solver chooses its own internal step size to keep
the estimated local error under the tolerance you asked for. When it
reports \code{1289 solver steps} for a run you asked to sample at 400
points, those 1289 are its own steps; the 400 are just where it stopped
to tell you the answer. \textbf{The output cadence does not affect the
physics.} That property is tested, not assumed.

\subsection{What the 7.8.0 upgrade brought}

This repository is \code{once-ere/rustSolveIt} with its 7.7.0 engine
replaced by the 7.8.0 one. Two things came with it.

\textbf{Four more solver families.}

\begin{center}
\begin{tabular}{lll}
\textbf{crate} & \textbf{solves} & \textbf{used here?}\\\hline
\code{cvode\_rs}  & ordinary differential equations (Adams / BDF) & yes\\
\code{arkode\_rs} & Runge--Kutta: explicit, implicit, IMEX, multirate, symplectic & yes\\
\code{ida\_rs}    & differential-algebraic systems $F(t,y,\dot y)=0$ & yes --- \kw{constrain} + \kw{method ida}\\
\code{kinsol\_rs} & nonlinear algebraic systems, Anderson acceleration & yes --- \kw{equilibrium}\\
\code{cvodes\_rs} & CVODE plus forward and adjoint sensitivity analysis & yes --- \kw{sensitivity}\\
\code{idas\_rs}   & IDA plus sensitivity analysis & yes --- \kw{sensitivity} while constrained\\
\end{tabular}
\end{center}

All six are compared byte-for-byte against the output of the original C
programs.

\subsection{The three questions that are not ``what happens next''}

\textbf{Four joints.} \kw{constrain} \emph{a b} is a rod (1 row, 5
freedoms left); \kw{ball} \emph{a b} shares a point (3 rows);
\kw{universal} \emph{a b u w} shares a point and keeps two shafts square,
a Cardan joint (4 rows); \kw{hinge} \emph{a b axis} shares a point
\emph{and} an axis, leaving exactly one freedom (5 rows) --- a door, a
knee, a pendulum. The last three grip orientation, which is why
\kw{method ida} carries the full 13-numbers-per-object state.

\textbf{A rigid rod, held exactly} --- \kw{constrain} \emph{a b} then
\kw{method ida}. A rod is a geometric fact, not a very stiff spring: the
bob is \emph{exactly} $L$ from the pivot, always. Saying so turns the
equations of motion from an ODE into a differential-algebraic equation,
which is what IDA solves. A body with \code{inverse\_mass = 0} becomes
an \textbf{anchor} --- it never moves and absorbs the reaction --- so
anchor + bob + rod is a pendulum. \kw{constraints} reports how far the
rod has strayed; over a 20-second swing it stays at roundoff.

\textbf{Where it comes to rest} --- \kw{equilibrium}. Not an
integration: KINSOL solves for a configuration where every free body has
zero net force, every anchor is where it started and every rod is the
right length, then puts the bodies there and stops them. It says nothing
about \emph{stability} --- a pencil on its point is an equilibrium too.

\textbf{How much the answer depends on an input} ---
\kw{sensitivity} 3 \code{"gravity.y"}. Running twice with slightly
different inputs and subtracting throws away most of your digits. This
integrates the derivative alongside the state instead. CVODES for an
ordinary system; IDAS when a \kw{constrain} is active.

\textbf{A more faithful API.} The 7.8.0 translation models C's opaque
pointers directly rather than smoothing them into idiomatic Rust. Vector
data is reached through a borrow guard, constructors return
\code{Option} where C returns a possibly-null pointer, and destructors
blank the handle the way \code{free(p); p = NULL;} does. None of that is
visible from the command language; the full call-by-call table is in
\code{PORT\_7.8.0\_PROVENANCE.md}.

\textbf{And the physics did not move.} All six self-checking physics
examples, all twelve collision scripts and all 59 dynamic notebooks
produce output that is byte-for-byte identical under 7.7.0 and 7.8.0.

\section{How collisions actually work}

\subsection{The problem with checking every frame}

The obvious way to detect a collision is to advance time a little, then
ask ``are these two objects overlapping?'' If yes, back up and resolve
it.

This fails in a specific, embarrassing way. Consider a bullet at
100\,m/s and a plate 5\,mm thick. In one hundredth of a second the
bullet moves a full metre --- two hundred times the plate's thickness.
Check before the step: not overlapping. Check after: not overlapping.
The bullet is now on the other side, and the program never noticed. This
is called \textbf{tunnelling}.

\subsection{What this program does instead}

It hands the solver a \textbf{continuous function} --- the signed gap
between two surfaces, positive when apart, negative when overlapping ---
and asks it to find the moment that function crosses zero. This is
\textbf{rootfinding}, a first-class feature of CVODE, not something
bolted on top.

Three things follow. \textbf{The impact time is exact}, not quantised to
a frame: the solver interpolates its own solution back to the crossing.
\textbf{The solver caps its own step size} while rootfinding is armed,
computed from the smallest feature among the paired bodies divided by
the fastest surface approach speed --- and it accounts for
\emph{acceleration}, so a body released from rest above a thin plate
still gets a finite cap. \textbf{Systems with nothing to collide are
unaffected}: with no collidable pairs the code path is bit-identical to
running with collisions off, and that is a tested property.

\subsection{What happens at the moment of contact}

The solver stops at the crossing and hands control back. The program
then reads the state at exactly that instant; computes the contact point
and normal from the two shapes' actual geometry (a \emph{support
function}, not a bounding sphere --- see example 13); applies an impulse
along the normal, sized by the two bodies' masses, inertias and
coefficients of restitution; and re-initialises the solver at the new
state. The impulse is applied at \emph{one shared point} for both
bodies, which is what makes angular momentum come out right for
off-centre hits.

\subsection{The Zeno problem}

A bouncing ball with restitution below 1 makes infinitely many bounces
in finite time. The program counts events per \emph{burst} --- a run of
events with no real flight between them. Past a threshold it forces the
remaining contacts plastic; past twice the threshold it projects out any
overlap and disarms rootfinding for the rest of the output interval. The
counting is per-burst rather than per-output-interval, because
per-interval counting made the physics depend on how often you asked for
output.

\section{Nineteen worked examples}\label{sec:examples}

Every transcript below is genuine program output. The scripts are in
\code{scripts/solveit/}:

\begin{lstlisting}
cargo run -p posim --release -- --script scripts/solveit/01_elastic_head_on.posim
\end{lstlisting}

\subsection{A head-on elastic collision between unequal masses}

Two spheres, masses 3 and 1, approach head-on at $+2$ and $-4$.
Elementary mechanics gives
$v_1' = ((m_1-m_2)u_1 + 2m_2u_2)/(m_1+m_2)$ and the mirror image for
$v_2'$. The answer is exact, and both energy and momentum must survive
the impulse untouched.

\begin{lstlisting}
In[5]:= energy
Out[5]= 14
In[6]:= momentum
Out[6]= [2, 0, 0]
In[7]:= run 2 steps 200
Out[7]= t = 2 (53 solver steps, 200 snapshots, |dE/E| = 0.000e0,
        1 collision(s))
In[8]:= get heavy.velocity
Out[8]= [-1, 0, 0]
In[9]:= get light.velocity
Out[9]= [5, 0, 0]
In[14]:= ((m1 - m2) * u1 + 2 * m2 * u2) / (m1 + m2)
Out[14]= -1
In[15]:= ((m2 - m1) * u2 + 2 * m1 * u1) / (m1 + m2)
Out[15]= 5
In[16]:= energy
Out[16]= 14
In[17]:= momentum
Out[17]= [2, 0, 0]
\end{lstlisting}

$-1$ and $5$ exactly, not approximately. Energy 14 before and 14 after;
momentum $[2,0,0]$ before and after. And $|dE/E| = 0$ --- the run lost
nothing at all. Note that \kw{let} makes a session variable and a bare
expression line is evaluated and printed, so the closed form is checked
\emph{inside the language}.

\subsection{Kepler's third law, four orbits at once}

For a circular orbit of radius $a$ about mass $M$, the period is
$T = 2\pi a^{3/2}/\sqrt{GM}$. Kepler's third law says $T^2/a^3$ is the
same number for every orbit around the same central mass. Each planet
is run for its own period and must come home.

\begin{lstlisting}
Out[5]=  [1.0000000044544244, 0, -0.00000014431704032681625]
Out[11]= [2.000000000309673,  0,  0.000000026045939587396316]
Out[17]= [3.0000000004230944, 0,  0.00000004071881288571222]
Out[23]= [4.000000000077065,  0,  0.00000007248745821777447]

In[24]:= 0.19869176531592203 * 0.19869176531592203 / 1
Out[24]= 0.039478417604357434
In[25]:= 0.5619851784832581 * 0.5619851784832581 / (2 * 2 * 2)
Out[25]= 0.039478417604357434
In[26]:= 1.0324326977181857 * 1.0324326977181857 / (3 * 3 * 3)
Out[26]= 0.039478417604357434
In[27]:= 1.5895341225273762 * 1.5895341225273762 / (4 * 4 * 4)
Out[27]= 0.039478417604357434
\end{lstlisting}

All four are the same 17 digits. Not ``close'' --- identical. That is
$4\pi^2/1000$.

\subsection{Bound, parabolic and hyperbolic, decided by a sign}

The shape of an orbit is decided by the specific orbital energy
$\varepsilon = v^2/2 - GM/r$: negative is an ellipse, exactly zero a
parabola, positive a hyperbola. The escape speed at $r$ is
$\sqrt{2GM/r}$; at $r=4$ with $GM=1000$ that is $\sqrt{500}$.

\begin{lstlisting}
In[7]:= 20 * 20 / 2 - gm / 4
Out[7]= -50
In[8]:= 22.360679774997898 * 22.360679774997898 / 2 - gm / 4
Out[8]= 0.00000000000002842170943040401
In[9]:= 24 * 24 / 2 - gm / 4
Out[9]= 38
\end{lstlisting}

$-50$, then $2.8\times10^{-14}$ (the closest a 64-bit float gets to zero
after squaring an irrational number), then $+38$. Three conics from one
formula.

\subsection{The restitution ladder, and honest energy loss}

A ball with restitution $e$ reaches $e^{2n}h$ after $n$ bounces. The
centre starts at 5.5 with radius 0.5, so it falls 5.0; $g=10$, $e=0.7$.
\textbf{\kw{run} takes a duration}, so the three runs are 1.7, 1.19,
0.833.

\begin{lstlisting}
In[7]:= get ball.position.y
Out[7]= 2.949999999869974
In[8]:= 5 * 0.7 * 0.7 + 0.5
Out[8]= 2.9499999999999997
In[10]:= get ball.position.y
Out[10]= 1.7004999998062575
In[11]:= 5 * 0.7 * 0.7 * 0.7 * 0.7 + 0.5
Out[11]= 1.7004999999999997
In[13]:= get ball.position.y
Out[13]= 1.0882449997748824
In[14]:= 5 * 0.7^6 + 0.5   (written out longhand in the script)
Out[14]= 1.0882449999999997
\end{lstlisting}

Three apexes, each matching to about $2\times10^{-10}$ --- and each
arrived at at the right \emph{moment}, which is the harder half.

Also notice $|dE/E| = 4.636\text{e-}1$. Energy is \textbf{not} conserved
here, and should not be: $e=0.7$ means each bounce deliberately throws
away 51\,\% of the kinetic energy. A conservation check is only
meaningful when the physics is conservative. This is the first place
most people misread the diagnostic.

\subsection{Cyclotron motion, and when to reach for BDF}

A charge $q$ moving at speed $v$ perpendicular to $B$ travels in a
circle of radius $r = m|v|/(|q|B)$ with period $T = 2\pi m/(|q|B)$. The
Lorentz force $q\,v\times B$ depends on velocity --- which rules out
symplectic methods entirely --- and is fast compared with anything else
in the problem. That is the textbook definition of a \emph{stiff}
system, and what \kw{method bdf} is for.

Half a period puts the ion diametrically opposite, so the distance from
the origin must be exactly $2r$.

\begin{lstlisting}
In[6]:= get ion.position
Out[6]= [-0.00000000020913167204927863, 1.1111111086955652, 0]
In[7]:= norm(ion.position)
Out[7]= 1.1111111086955652
In[8]:= 2 * (2 * 5 / (3 * 6))
Out[8]= 1.1111111111111112
In[10]:= get ion.position          (after the second half period)
Out[10]= [0.000000000408039691444928, 0.000000005381884658509583, 0]
In[11]:= get ion.velocity
Out[11]= [4.999999951563042, 0.0000000036723595761567474, 0]
\end{lstlisting}

Right to 8 significant figures; after a full period the ion is back at
the origin to 5 parts in a billion with its original velocity restored.
The magnetic force does no work, and the speed confirms it.

\subsection{Symplectic versus adaptive, over a long run}

CVODE Adams controls the \emph{local} error at every step. A
\textbf{symplectic} method controls something else entirely: it exactly
preserves the geometric structure of Hamiltonian mechanics, so the
energy error does not accumulate --- it oscillates around a constant.

\begin{lstlisting}
In[5]:= run 60 steps 60                         (METHOD ADAMS)
Out[5]= t = 60 (507075 solver steps, 60 snapshots, |dE/E| = 1.459e-8)
In[6]:= energy
Out[6]= -0.20000000291885495

In[12]:= run 60 steps 60      (METHOD SPRK MCLACHLAN_4_4 0.001)
Out[12]= t = 60 (60000 solver steps, 60 snapshots, |dE/E| = 8.957e-11)
In[13]:= energy
Out[13]= -0.20000000001790574
\end{lstlisting}

The symplectic method took \textbf{60{,}000} steps to the adaptive
method's \textbf{507{,}075}, and its energy error is \textbf{163 times
smaller}. Eight times less work, two orders of magnitude better answer.

\textbf{The catch, and it is a real one.} SPRK needs a \emph{separable}
Hamiltonian. A magnetic field breaks that, and so does rotation. The
program does not let you get this wrong quietly:

\begin{lstlisting}
Err: SPRK method requires a separable Hamiltonian: magnetic field B must
     be zero (the Lorentz force q v x B is velocity-dependent);
     use METHOD ADAMS or BDF
\end{lstlisting}

The error names the feature that blocks it. That refusal is why the
previous example uses BDF.

\subsection{The tennis-racket theorem, measured}

A rigid body has three principal axes with three different moments of
inertia. Spin it about the largest or the smallest and the motion is
stable. Spin it about the \textbf{middle} one and it is not: the body
flips end over end, over and over. Cosmonaut Vladimir Dzhanibekov
noticed this on Salyut 7 in 1985 with a wing nut.

The flip is a genuine instability, so it is exquisitely sensitive to
integration error --- and yet the angular momentum must be
\emph{exactly} conserved throughout, because nothing is applying a
torque. Those two demands pull in opposite directions.

\begin{lstlisting}
In[2]:= angmom
Out[2]= [0.02, 3, 0.02]
In[3]:= energy
Out[3]= 5.148625491836798

racket.orientation.x, sampled every 0.5 from t = 0.5 to t = 4.0:
   0.013359726984815766
   0.009214316848229325
  -0.06928171314906612
  -0.32837233114496406
  -0.4991829732950659
   0.1093367245470612
   0.8230403847701171
   0.9661018228588313

In[20]:= angmom
Out[20]= [0.02, 3, 0.02]
In[21]:= energy
Out[21]= 5.148625488700016
\end{lstlisting}

The $x$ component of the orientation quaternion --- near enough, ``how
far over has it rolled'' --- sweeps from 0.013 through negative
territory to 0.966. That is the flip. And through all of it,
\kw{angmom} reads $[0.02, 3, 0.02]$ at the end exactly as it did at the
start: not to eight decimals, the identical printed value.

You can watch this one: \code{videos/tumbling\_racket.html}.

\subsection{When energy \emph{should} grow}

A magnetised body in a field feels a torque $\tau = (RMR^{T})B$. The
field is an outside agent: it does work, so kinetic energy grows and
angular momentum about the origin is \emph{not} conserved. Half of
understanding a conservation diagnostic is knowing when it should fail.

\begin{lstlisting}
In[3]:= new cuboid as rotor { mass = 1, half_extents = [1, 0.4, 0.4],
          magnetic_moment_tensor = [[0, 0, 0.5], [0, 0, 0], [-0.5, 0, 0]] }
In[4]:= energy
Out[4]= 0
In[5]:= angmom
Out[5]= [0, 0, 0]
In[6]:= run 3 steps 30
Out[6]= t = 3 (550 solver steps, 30 snapshots, |dE/E| = 9.928e-1)
In[7]:= energy
Out[7]= 0.9928369509454675
In[8]:= angmom
Out[8]= [0.46022300703213415, 0, 0]
\end{lstlisting}

Starts at rest, ends spinning. The initial torque is
$M\!\cdot\!B = (0.5\cdot2, 0, 0) = (1,0,0)$, and the angular momentum
duly grows along $x$.

\textbf{The trap.} With
\code{[[0, 0.5, 0], [-0.5, 0, 0], [0, 0, 0]]} --- a perfectly
reasonable-looking antisymmetric tensor --- the \emph{third column} is
zero, $M\!\cdot\!B$ is the zero vector, and absolutely nothing happens.
No error, no warning, just a body that sits there. Multiply the tensor
by your field by hand before you blame the solver.

\subsection{Newton's cradle}

Five identical spheres, four of them touching. Roll one into the line
and \emph{only the far one} leaves, at exactly the incoming speed.
Momentum alone does not force this --- two balls leaving at half speed
would conserve momentum fine. It is momentum \textbf{and} energy
together that pick out the answer.

\begin{lstlisting}
In[8]:= energy
Out[8]= 0.5
In[9]:= momentum
Out[9]= [1, 0, 0]
In[10]:= run 6 steps 600
Out[10]= t = 6 (30 solver steps, 600 snapshots, |dE/E| = 0.000e0,
         4 collision(s))
In[11..14]:= get b0..b3 .velocity
Out[11..14]= [0, 0, 0]   (all four)
In[15]:= get b4.velocity
Out[15]= [1, 0, 0]
In[16]:= energy
Out[16]= 0.5
In[17]:= momentum
Out[17]= [1, 0, 0]
\end{lstlisting}

Four impulses, and not one bit lost.

\subsection{The bullet that does not tunnel}

A 4\,mm bullet at 100\,m/s meets a 5\,mm plate. Between output frames
the bullet travels a metre --- two hundred plate thicknesses. This is
precisely the case that defeats frame-sampled collision detection.

\begin{lstlisting}
In[5]:= run 0.1 steps 10
Out[5]= t = 0.1 (10003 solver steps, 10 snapshots, |dE/E| = 0.000e0,
        1 collision(s))
In[6]:= get bullet.position.y
Out[6]= 8.009000000001455
In[7]:= get bullet.velocity.y
Out[7]= 100
In[9]:= contacts
Out[9]= contact0: obj0 <-> obj1 at t = 0.01995500000000067
  point  = [0, 0.0025, 0]
  normal = [-0, 1, -0]  (from obj0 toward obj1)
  approach speed = 100, impulse = 2
\end{lstlisting}

The bullet had 1.9955\,m of gap to close at 100\,m/s, so the exact
answer is 0.019955. Look also at the step count: \textbf{10{,}003 solver
steps} for ten output points. That is the anti-tunnelling cap at work
--- the solver was told it may not take a step large enough to skip the
plate, and it obeyed.

\subsection{Lagrange's equilateral three-body solution}

Three equal masses at the corners of an equilateral triangle, each
moving at the right speed, rotate rigidly forever --- one of only a
handful of exactly solvable configurations of the three-body problem,
found by Lagrange in 1772. The same geometry is why Jupiter has Trojan
asteroids. With $G=m=1$ and side $a$, $\omega = \sqrt{3Gm/a^3}$.

There are three independent things to check: each body comes home, the
centre of mass never moves, and the triangle stays equilateral.

\begin{lstlisting}
In[6]:= energy
Out[6]= -1.4999999999984994
In[7]:= run 3.6275987284684357 steps 400
Out[7]= t = 3.6275987284684357 (1289 solver steps, 400 snapshots,
        |dE/E| = 4.689e-10)
In[8]:=  get a.position
Out[8]=  [0.5773502692574497, 0, -0.0000000025627822854319693]
In[9]:=  get b.position
Out[9]=  [-0.2886751324104037, 0, 0.5000000013400956]
In[10]:= get c.position
Out[10]= [-0.28867513684704826, 0, -0.49999999877731977]
In[11]:= com
Out[11]= [-7.586524001605236e-16, 0, -2.146431180941969e-15]
In[12]:= norm(a.position - b.position)
Out[12]= 1.0000000001184224
\end{lstlisting}

After a full revolution all three are back within about
$3\times10^{-9}$. The centre of mass has moved by $2\times10^{-15}$ ---
five million times less, which is what you would expect, since it is
protected by momentum conservation rather than by accuracy. And the side
length is still 1 to 1 part in $10^{10}$.

Note also \code{norm(a.position - b.position)}: you can do vector
algebra on live simulator fields directly in the language.

\subsection{A user-defined constructor, and an inertia tensor}

A \kw{dumbbell} is two solid spheres joined by a rod, treated as
\textbf{one rigid body}. Its inertia about the centre of mass is the
exact composite: $\tfrac25 mr^2$ for each sphere, plus $mz^2$ by the
parallel-axis theorem, plus the rod's own terms. The inertia tensor is
where compound-body code usually goes wrong, and the answer is a hand
calculation.

\begin{lstlisting}
In[1]:= def barbell(m1 = 1, m2 = 1, m_rod = 0.5, r1 = 0.3, r2 = 0.3,
                    rod_r = 0.05, len = 2) {
  new dumbbell as bar { m1 = m1, m2 = m2, m_rod = m_rod, r1 = r1,
                        r2 = r2, rod_radius = rod_r, length = len }
}
Out[1]= function barbell(7 parameter(s)) defined - 1 body line(s)
In[3]:= barbell(2, 2, 0, 0.3, 0.3, 0.05, 2)
Out[3]= obj0 as bar
In[4]:= list
Out[4]= obj0: dumbbell r1=0.3 r2=0.3 rod_r=0.05 len=2, mass=4,
        charge=0, pos=[0, 0, 0]
In[9]:= get bar.inertia_tensor
Out[9]= [[4.144, 0, 0], [0, 4.144, 0], [0, 0, 0.144]]
In[10]:= 2 * (0.4 * 2 * 0.3 * 0.3 + 2 * 1 * 1)
Out[10]= 4.144
In[11]:= 2 * (0.4 * 2 * 0.3 * 0.3)
Out[11]= 0.144
\end{lstlisting}

$4.144$ and $0.144$, exactly. The off-diagonal entries are exactly zero,
as they must be for a body symmetric about its $z$ axis. Every round
shape in this program --- torus, disk, cylinder, dumbbell --- is
symmetric about its \textbf{local $z$ axis}; that is the convention.

\subsection{Geometry a bounding sphere gets wrong}

An axis-aligned torus with outer radius 2 exactly inscribes a cube of
inner side 4. Tilt that same torus so its axis points along
$(1,1,1)/\sqrt3$ and its extent along each coordinate axis \emph{drops}
to $1.5\sqrt{2/3} + 0.5 \approx 1.7247$, well inside the walls.

A bounding sphere would say the opposite: the torus's bounding sphere
has radius 2 no matter how you turn it. Getting this right requires the
actual \textbf{support function} of the shape --- the farthest point of
the body in a given direction --- evaluated in the rotated frame.

\begin{lstlisting}
In[2]:= box 4
Out[2]= box: inner size 4 x 4 x 4 - six static walls obj0..obj5 with
        inverse_mass = 0
In[7]:= new torus as tilt { mass = 1, ring_radius = 1.5,
          tube_radius = 0.5, position = [0, 0, 0],
          orientation = [0.8880738339771154, -0.32505758367186816,
                         0.32505758367186816, 0] }
In[8]:= run 1 steps 10
Out[8]= t = 1 (4 solver steps, 10 snapshots, |dE/E| = 0.000e0)
In[10]:= get system.collisions
Out[10]= 0
In[11]:= 1.5 * sqrt(2 / 3) + 0.5
Out[11]= 1.724744871391589
\end{lstlisting}

Zero collisions. The clearance ($2 - 1.7247 = 0.2753$ per side) is
exactly what the geometry predicts. Note also the \kw{box} machinery:
six wall slabs, created automatically, with \code{inverse\_mass = 0} ---
\emph{infinite} mass, so they absorb any impulse without ever moving.

\subsection{The infinite square well, solved in the language}

A particle confined to a box of width $L$ has energies
$E_n = n^2\pi^2\hbar^2/(2mL^2)$. \kw{qm grid} pins the wavefunction to
zero just outside the domain, which \emph{is} an infinite square well
--- the program is solving the real problem, not a special case of it.

\begin{lstlisting}
In[1]:= qm grid 0 1 800
In[2]:= qm potential zero
In[5]:= qm states 4
Out[5]= 4 lowest bound state(s):
  E[0] =  4.934795874467
  E[1] = 19.739107587630
  E[2] = 44.412707408131
  E[3] = 78.955215788088
In[6..9]:= n^2 * pi * pi / 2  for n = 1..4
Out[6]=  4.934802200544679
Out[7]= 19.739208802178716
Out[8]= 44.41321980490211
Out[9]= 78.95683520871486
\end{lstlisting}

The agreement is to about 1.3 parts in a million, and it gets
\emph{worse} for higher levels. That is not sloppiness: it is the
finite-difference approximation to the second derivative, whose error
grows with the curvature of the state, and which shrinks as $h^2$ when
you refine the grid.

A physical cross-check on the loaded state:

\begin{lstlisting}
In[11]:= qm norm
Out[11]= 1.000000000000345
In[12]:= qm position
Out[12]= 0.500000000001550
In[14]:= qm prob 0 0.5
Out[14]= 0.499999999998032
\end{lstlisting}

Total probability 1, mean position dead centre, exactly half the
probability in each half of the box --- all to 12 digits. The
\emph{eigenvalue} carries discretisation error; the \emph{state} is
properly normalised and properly symmetric.

\subsection{Tunnelling, and watching a grid converge}

A particle with energy $E$ below a barrier of height $V_0$ and width $a$
still gets through, with
\[
T = \frac{1}{1 + \dfrac{V_0^2\sinh^2(\kappa a)}{4E(V_0-E)}},
\qquad \kappa = \frac{\sqrt{2m(V_0-E)}}{\hbar}.
\]
With $V_0=4$, $a=1$, $E=2$, $\hbar=m=1$ that is $0.07065082485\ldots$

Transmission depends on the barrier width \emph{exponentially}, so a
grid that gets the width slightly wrong gets $T$ badly wrong. That makes
it an unusually sharp probe of the discretisation.

\begin{lstlisting}
In[7]:= 1 / (1 + v0 * v0 * s * s / (4 * e * (v0 - e)))
Out[7]= 0.07065082485316446

qm grid -30 30  4000 -> T = 0.069368404960, T + R = 1.000000000000
qm grid -30 30 16000 -> T = 0.070328037206, T + R = 1.000000000000
qm grid -30 30 64000 -> T = 0.070569990790, T + R = 1.000000000000
\end{lstlisting}

\begin{center}
\begin{tabular}{rrr}
\textbf{points} & \textbf{$T$} & \textbf{error}\\\hline
4{,}000  & 0.069368 & $1.3\times10^{-3}$\\
16{,}000 & 0.070328 & $3.2\times10^{-4}$\\
64{,}000 & 0.070570 & $8.1\times10^{-5}$\\
\end{tabular}
\end{center}

The error falls by a factor of $\sim4$ each time the grid is refined by
a factor of 4 --- first-order convergence, exactly what a
staircase-sampled barrier gives you. \textbf{This is what
``discretisation error'' looks like from the outside}, and it is how you
tell it from a bug: a bug does not converge.

Notice too that $T + R = 1.000000000000$ on every grid. Probability is
conserved exactly regardless of how coarse the grid is, because that is
a structural property of the transfer matrix, not an accuracy property.
The two kinds of correctness are independent.

Above the barrier there are resonant energies where it becomes
perfectly transparent:

\begin{lstlisting}
In[19]:= qm scan 4.1 20 200
Out[19]= 200 energies in [4.1, 20] (0 refused), T from 3.535e-1 to 1.000e0
  1 resonance(s), strongest first:
    E = 8.893969849   T = 0.999992036
\end{lstlisting}

\subsection{The special-function library, checked by identities}

The way to test such a library is not to compare against a table (tables
have typos) but against \textbf{identities} --- relations a wrong
implementation cannot satisfy by accident.

\begin{lstlisting}
In[1]:= bessel_j(0, 2.404825557695773)
Out[1]= 0
In[2]:= bessel_j_nu(1, 3) * bessel_y_nu(0, 3)
        - bessel_j_nu(0, 3) * bessel_y_nu(1, 3)
Out[2]= 0.21220659078919374 + 0i
In[3]:= 2 / (pi * 3)
Out[3]= 0.2122065907891938
In[4]:= legendre_p(5, 1)
Out[4]= 1
In[5]:= legendre_p(5, 0.906179845938664)
Out[5]= -1.7763568394002506e-16
In[6]:= gamma_z(0.5) * gamma_z(0.5)
Out[6]= 3.1415926535898397 + 0i
In[7]:= pi
Out[7]= 3.141592653589793
In[8]:= sph_harm_real(1, 0, 0, 0)
Out[8]= 0.48860251190291987
In[9]:= sqrt(3 / (4 * pi))
Out[9]= 0.4886025119029199
In[10]:= clebsch_gordan(0.5, 0.5, 0.5, -0.5, 1, 0)
Out[10]= 0.7071067811865475
In[11]:= sqrt(0.5)
Out[11]= 0.7071067811865476
In[12]:= gauss_legendre(4)
Out[12]= [[-0.8611363115940526, -0.3399810435848563,
            0.3399810435848563,  0.8611363115940526],
          [ 0.34785484513745374, 0.6521451548625461,
            0.6521451548625461,  0.34785484513745374]]
\end{lstlisting}

Line by line: $J_0$ at its first zero is a clean 0. The Wronskian
identity $J_{\nu+1}Y_\nu - J_\nu Y_{\nu+1} = 2/(\pi x)$ matches to 16
digits --- a genuinely strong test, since it couples four different
function evaluations. $P_5(1)=1$ exactly, and $0.906179845938664$ really
is one of its roots (residual one bit). $\Gamma(\tfrac12)^2 = \pi$ to 14
digits, via the \emph{complex} gamma evaluated on the real axis. The
spherical harmonic $Y_1^0$ at $\theta=0$ is $\sqrt{3/4\pi}$ to 15
digits. Two spin-$\tfrac12$ particles coupling to the $m=0$ triplet have
Clebsch--Gordan coefficient $1/\sqrt2$ to the last bit. And the 4-point
Gauss--Legendre rule reproduces the classical nodes and weights, whose
sum is 2 --- the length of $[-1,1]$.

\textbf{One design decision worth knowing.} Where an argument is an
integer order, a fractional value is refused rather than truncated:

\begin{lstlisting}
In[1]:= hermite_h(2.5, 1)
Err[1]: hermite_h(): argument 1 must be a whole number
        (an integer order), got 2.5
\end{lstlisting}

Truncating would have returned a confident, wrong number with no way for
you to notice. Angular momenta, on the other hand, \emph{may} be
half-integers, so the Wigner and Clebsch--Gordan routines take plain
numbers.

\subsection{A pendulum as a constraint, not a force}

The small-amplitude period is $T = 2\pi\sqrt{L/g}$; released 0.02\,rad
off vertical with $L=1$, $g=9.81$ that is $2.0060666807106475$\,s, and
after exactly one period the bob must be back where it started. The rod
is not a force but a geometric condition, and the only way to hold it
exactly is to solve the DAE.

\begin{lstlisting}
In[4]:= new sphere as pivot { mass = 1, radius = 0.02,
                              position = [0, 0, 0], inverse_mass = 0 }
In[5]:= new sphere as bob   { mass = 1, radius = 0.05,
          position = [0.019998666693333084, -0.9998000066665778, 0] }
In[6]:= constrain pivot bob
Out[6]= constraint0: obj0 <-> obj1 held at 1
        (METHOD IDA is required to integrate it)
In[8]:= method ida
Out[8]= method = IDA (constrained DAE, GGL index-2)
In[9]:= run 2.0060666807106475 steps 200
Out[9]= t = 2.0060666807106475 (339 solver steps, 200 snapshots,
        |dE/E| = 2.530e-12)
In[10]:= get bob.position
Out[10]= [0.019998666320588818, -0.9998000066740419, 0]
In[11]:= constraints
Out[11]= constraint0: obj0 <-> obj1 length 1
         (currently 1.0000000000000082)
\end{lstlisting}

The bob came back to within $3.7\times10^{-10}$ after a full period, and
the rod reads $1.0000000000000082$ --- eight parts in $10^{15}$, one
bit, after 339 solver steps. That is not luck: the formulation carries
\emph{both} the rod's length and its rate of change as equations. The
cheaper scheme that constrains only the acceleration lets the length
drift quadratically, and you find out much later.

The \code{inverse\_mass = 0} on the pivot is what makes it a pendulum
rather than two bodies tumbling about their shared centre of mass joined
by a stick.

\subsection{Where it rests, and what the answer depends on}

A pendulum released anywhere comes to rest hanging straight down, one
rod-length below the pivot:

\begin{lstlisting}
In[7]:= equilibrium
Out[7]= equilibrium found in 17 Newton iteration(s), 67 residual
        evaluation(s); largest net force on any free body =
        7.459152323898993e-13, worst |g| = 1.9317880628477724e-14
In[8]:= get bob.position
Out[8]= [0.0000000000000735262479725006, -0.9999999999999807, 0]
\end{lstlisting}

Released $57^\circ$ off vertical, the bob lands straight down to 13
digits. Now the derivative, on a free body, where
$\partial y/\partial g = T^2/2 = 4.5$ at $T=3$:

\begin{lstlisting}
In[15]:= sensitivity 3 "gravity.y" "mass 0"
Out[15]= t = 3 (CVODES, 129 solver steps)
d/d(gravity.y):
  obj0 position [0, 4.500000056696235, 0]
d/d(mass 0):
  obj0 position [0, 0, 0]
In[16]:= get stone.position
Out[16]= [3.000000000000001, -44.14500000000045, 0]
\end{lstlisting}

$4.500000056696235$ against an analytic $4.5$ --- 1.3 parts in $10^8$,
carried alongside a trajectory that itself landed on
$-44.14500000000045$ where $-\tfrac12\cdot 9.81\cdot 9 = -44.145$.

And the second derivative is \textbf{exactly zero}. In uniform gravity
every mass accelerates equally, so the trajectory does not depend on the
mass at all. A finite-difference estimate would have returned some small
number and left you guessing whether it was noise or physics. The
sensitivity equations return the zero.

\subsection{A door on a hinge}

A hinge fixes a point \emph{and} an axis, leaving one freedom. A hinged
rigid body is a \textbf{compound} pendulum: its small-amplitude period is
$T = 2\pi\sqrt{I_{\text{pivot}}/(mgd)}$ with
$I_{\text{pivot}} = I_{\text{com}} + md^2$ --- the moment of inertia
about the pivot, not just the distance to the centre of mass. The
point-mass formula is wrong by 15\,\% for the slab below.

\begin{lstlisting}
In[4]:= new sphere as jamb { mass = 1, radius = 0.02,
                             position = [0, 0, 0], inverse_mass = 0 }
In[5]:= new cuboid as door { mass = 1, half_extents = [0.2, 0.4, 0.2],
          position = [0.0199986666933331, -0.9998000066665778, 0] }
In[6]:= hinge jamb door [0, 0, 1]
In[8]:= method ida
Out[8]= method = IDA (constrained DAE, GGL index-2)
In[9]:= run 1 steps 10
Out[9]= t = 1 (70 solver steps, 10 snapshots, |dE/E| = 1.613e-9)
In[10]:= constraints
Out[10]= constraint0: hinge obj0 <-> obj1, 5 row(s)
worst |g| = 2.73750133672479e-10, worst |g_dot| = 2.3769240437118873e-9
In[11]:= get door.angular_momentum
Out[11]= [0, 0, 0.003742219622967182]
\end{lstlisting}

The angular momentum is $[0, 0, 0.0037]$ --- \emph{exactly} zero on $x$
and $y$. The two rows a hinge has over a ball joint are the ones that
forbid those axes, and they are doing their job to the last bit. The
joint is held to $2.7\times10^{-10}$ after 70 solver steps, and energy to
1.6 parts in $10^{9}$; run one compound period and the slab returns to
within $3\times10^{-8}$.

The \code{inverse\_mass = 0} on the jamb is what makes it a door: an
anchor, immovable, absorbing whatever the hinge pulls with.

\textbf{A joint constrains velocity, not just position.} A ball joint
says the two bodies share a point, so at the velocity level
$v_i + \omega_i\times r_i = v_j + \omega_j\times r_j$: a body turning
about a pivot offset from its centre \emph{must have its centre moving}.
Give it a spin and leave its velocity at zero and the state is not on the
constraint manifold. The run projects the starting velocities onto it ---
the smallest mass-weighted change that satisfies the joint, which is
precisely the impulse a real coupling delivers when clutched onto a
spinning shaft --- and reports how big that change was. For a cube spun
at 3\,rad/s on a half-metre arm the projection leaves the turn nearly
untouched and sets the \emph{centre} moving instead: the pivot was
running at 1.5\,m/s, and giving a 1\,kg body some velocity is cheaper
than fighting the spin.

\textbf{One real limit.} Orientation joints carry a tolerance floor of
$\texttt{rtol}=10^{-6}$: the index-2 system has an accuracy ceiling, and
across twelve compound pendulums the boundary is sharp --- $10^{-6}$
converges in every one, $10^{-8}$ in none.

\section{Watching it: the scene window and browser videos}

\subsection{The live window}

Type \kw{scene create} and a browser tab opens. That tab \emph{is} the
simulator's window: it draws every body on a 3-D canvas, with a toolbar
(Start, Pause, Stop, Reverse, Reset, single-step, a \emph{dt} box, zoom,
grid/trails/labels toggles) and a status bar showing the playback mode,
the time, \emph{dt}, the total energy, the body count, the camera and
the frame rate.

Behind it is a web server written on the Rust standard library ---
including the HTTP parsing, the WebSocket handshake, the SHA-1 and the
base64. No dependency, and the page itself is vanilla JavaScript and
canvas with no CDN fetches.

Two things about it are worth knowing. \textbf{The window evolves its
own copy of the system}: typing \kw{step} in the notebook does not move
the window, and pressing Start in the window does not move the notebook.
That isolation is what lets the window run without ever locking the
language. And \textbf{Reverse is replay, not negative-time
integration}: the playback keeps a ring of snapshots and walks it
backward, returning to $t=0$ bit-identically.

\subsection{Recorded videos}

A live window needs a live program. For something you can send to
someone else, record it:

\begin{lstlisting}
cargo build --release -p posim
recorder/src/record_video.py videos/scenes/kepler_ellipse.posim \
     -o videos/kepler_ellipse.html --frames 360 --dt 0.02 \
     --title "Kepler orbit, e = 0.6"
\end{lstlisting}

The result is a single HTML file: the frames embedded as data plus a
canvas player. It fetches nothing, so it works from \code{file://} on a
machine with no network, forever. The recorder drives
\code{posim -{}-machine} and asks it to \kw{step} --- \textbf{every
advance is a real SUNDIALS step}. The tool is a camera, not a physics
engine.

\begin{center}
\begin{tabular}{p{3.6cm}p{6.0cm}p{4.6cm}}
\textbf{file} & \textbf{what to watch} & \textbf{measured}\\\hline
\code{kepler\_ellipse.html} & the speed swinging between perihelion and
aphelion & $|dE|/E = 9.8\times10^{-8}$,
$|dL|/|L| = 1.3\times10^{-7}$\\
\code{tumbling\_racket.html} & the tennis-racket flip, seen from outside
& $|d|L||/|L| = 0$ \textbf{exactly}; $|dE|/E = 6.4\times10^{-9}$\\
\code{box\_of\_shapes.html} & a cylinder, a disk and a cuboid in a rigid
\kw{box} 4; gold arrows are the analytic contact normals, sized by
impulse & 36 collisions, $|dE|/E = 3.4\times10^{-16}$\\
\code{double\_pendulum\_hinges.html} & \textbf{two \kw{hinge} joints}
assembled into the classic chaotic linkage --- gold rings mark the
joints, gold lines their axes & the joints hold to
$|g| = 5.6\times10^{-8}$ through 400 IDA steps\\
\code{universal\_joint.html} & a \textbf{\kw{universal} joint} carrying
a driven shaft's rotation across to a second shaft, braced by a rod to
a post & the bend stops at $\cos\beta = 0.6000004$ against a geometric
bound of exactly $0.6$; three joints hold to $|g| = 4.0\times10^{-7}$\\
\code{ball\_joint\_chain.html} & four links on \textbf{\kw{ball} joints}
--- the chain whirls out of the plane it started on, which a hinged
chain cannot do & the four joints hold to $|g| = 3.3\times10^{-9}$;
$|z|$ runs from exactly $0$ to $1.7147$\\
\code{rod\_pendulum\_chain.html} & four bobs on four
\textbf{\kw{constrain} rods} --- one row each, the cheapest linkage
in the language, going chaotic & run continuously at the default
tolerance the rods hold to $|g| = 5.4\times10^{-15}$, i.e. roundoff\\
\code{spinning\_top.html} & a top on a \textbf{\kw{ball} joint},
precessing under gravity & $1.020440$ rad/s against the exact
$\Omega = Mgr/(I_3\omega_3) = 1.020408$, 3 parts in $10^{5}$\\
\code{gyroscope\_gimbal.html} & a rotor in \textbf{two gimbal
rings}, three perpendicular \kw{hinge} axes & total $L\cdot\hat y$
conserved to $1.4\times10^{-14}$; every centre stays on the pivot to
$1.2\times10^{-34}$\\
\code{cardan\_gear.html} & \textbf{Cardan gears} --- a wheel in a
ring of twice its radius, rolling on a \kw{gear} row, its rim point
tracing a straight line & line held to $1.1\times10^{-8}$; stroke
exactly $\pm 2r$\\
\code{rack\_and\_pinion.html} & a weight on a \textbf{\kw{rack}}
winding up a flywheel, on a \textbf{\kw{prismatic}} guide & falls at
exactly $g/2$, and at the same rate for two different pitch radii\\
\code{piston\_crankshaft.html} & the \textbf{slider-crank} --- a
piston, a connecting rod and a crankshaft & follows the exact
kinematics to $8.4\times10^{-8}$; stroke exactly $L-a$ to $L+a$\\

\code{cardan\_compass.html} & the same rings with a
\textbf{pendulous} bowl --- a ship's compass, seeking level &
periods $1.878587$ and $2.307339$ s against measured $1.883426$ and
$2.313653$\\
\end{tabular}
\end{center}

That last one is the one to open if you only open one, because it shows
a joint doing something a rod cannot. A \kw{universal} holds a shared
point and one right angle between its trunnions; it deliberately does
\emph{not} hold the two shafts straight. Left dangling, the output
shaft therefore does not transmit rotation at all --- it swings down
past the joint like a pendulum and folds back at $176^\circ$. Bracing
it with a second bearing looks right and fails at once, because
$\kw{hinge}\,5 + \kw{universal}\,4 + \kw{hinge}\,5$ is 14 rows on 12
freedoms and three of them are redundant, which leaves the constrained
solve singular. One \kw{constrain} --- a single row --- braces it
instead, and pins the bend to a closed form: the shaft sweeps a cone,
and $\cos\beta$ cannot go below $0.6$. The recording touches
$0.6000004$. Grammar \S12.10 works the arithmetic through.

The player has Play/Pause (or Space), frame stepping (or
$\leftarrow\rightarrow$), a scrub bar, speed from $0.25\times$ to
$4\times$, drag to orbit, wheel to zoom, and toggles for trails, labels
and contact arrows. The corner readout shows the frame number, $t$, $E$,
$|P|$, $|L|$ and the collision count \emph{for the frame you are on} ---
so you can stop on the moment something looks wrong and read the
conserved quantities off it.

\section{How you know the numbers are right}

\subsection{The solver library against its C original}

\code{sundials\_rs/} ships the upstream SUNDIALS example programs,
translated. Each one is run and its output \textbf{diffed byte-for-byte}
against the output of the original C program. The results, the
divergences, and the reasons for each divergence are in
\code{sundials\_rs/VERIFICATION.md}, \code{sundials\_rs/differences/}
and \code{sundials\_rs/c-results/}.

Two consequences of taking this seriously: printed floating-point values
go through helpers that reproduce C's \code{printf} conversions exactly
(Rust's own \code{\{:e\}} formats exponents differently), and
\code{pow} was made host-independent so that at least that function
cannot vary between machines.

\subsection{The simulator against closed-form physics}

Six example programs, each printing \code{SUCCESS} or \code{FAILURE} and
exiting nonzero if it fails: \code{kepler\_orbit},
\code{outer\_solar\_system}, \code{tumbling\_body},
\code{charged\_in\_b\_field}, \code{newtons\_cradle},
\code{bouncing\_ball\_restitution}.

\subsection{The test suite}

605 tests: 40 library, 19 collision, 9 conservation, 109 language, 92
quantum, 233 special-function, 11 vendored identities and 55
documentation examples that are compiled and run as written. The
collision tests in particular are analytic rather than
regression-style: they assert the time of impact against a closed form,
not against last week's output.

\subsection{The notebooks}

\code{dynamic\_notebooks/} holds 59 runnable sessions, including 34
problems from Routh's \emph{Dynamics of a Particle} (1898) and
\emph{Dynamics of a System of Rigid Bodies} (1905). Each derives its
closed-form answer in its own header and then checks the integrator
against it --- so the numbers in the catalogue are measured, not quoted.

\subsection{And for the 7.8.0 upgrade specifically}

Everything above was run against \textbf{both} the old 7.7.0 engine and
the new 7.8.0 one, and the outputs diffed. They are identical byte for
byte. See \code{PORT\_7.8.0\_PROVENANCE.md} and
\code{evidence/port-7.8.0/}.

\section{When something goes wrong}

Errors name the column, the field, or the feature at fault, and never
abort your session.

\begin{center}
\begin{tabular}{p{5.0cm}p{9.6cm}}
\textbf{symptom} & \textbf{what is actually happening}\\\hline
$|dE/E|$ is huge and you expected conservation & Check
\code{system.g\_constant} (\S\ref{sec:gravity}), the restitution of your
bodies, and whether a field is doing work.\\
\kw{run} overshoots where you meant to stop & \kw{run} takes a
\textbf{duration}.\\
bodies pass through each other & You did not \kw{collide on}.\\
\code{SPRK method requires a separable Hamiltonian} & A magnetic field,
a magnetic moment, an external torque or a spinning body is present.
The message names which. Use \kw{method adams} or \kw{bdf}.\\
\code{no object obj7} after a \kw{del} & Deleting renumbers. Use
\kw{new} \dots\ \kw{as} \emph{name}.\\
\code{unknown name `x'} & A bare identifier resolves at execution time.
Bind it with \kw{let}, pass it as a function parameter, or use
\emph{name}\code{.}\emph{field}.\\
a magnetised body will not spin & Multiply your moment tensor by your
field by hand --- you probably have a zero column.\\
a quantum result is off by a fraction of a percent & Refine the grid and
see whether it converges. Discretisation error converges; a bug does
not.\\
\code{probability accumulating near the wall} & Your wavepacket has
reached the domain edge and is reflecting. Widen the grid or add
\kw{qm absorb}.\\
the scene window does not open & Set \code{POSIM\_NO\_BROWSER=1} if you
are headless, or open the printed \code{http://127.0.0.1:<port>/}
yourself.\\
\end{tabular}
\end{center}

Anything the language accepts is documented in \code{grammar.md} ---
every command, every field, every builtin, with the complete EBNF
grammar and a further eighteen worked examples.

\section{Where everything lives}

\begin{center}
\begin{tabular}{p{4.4cm}p{10.2cm}}
\textbf{path} & \textbf{what}\\\hline
\code{posim/} & the command language, the graphics window, the machine protocol\\
\code{physical\_object/} & the physics: the union type, collisions, all time integration\\
\code{quantum/} & 1-D, 2-D and 3-D quantum mechanics\\
\code{special\_functions/} & Bessel, Legendre, Hankel, Airy, gamma, Wigner, quadrature, eigenproblems\\
\code{sundials\_rs/} & the pure-Rust SUNDIALS 7.8.0 engine (vendored, read-only)\\
\code{jupyter/} & the JupyterLab kernel\\
\code{dynamic\_notebooks/} & 59 runnable sessions, incl.\ 34 Routh problems\\
\code{scripts/solveit/} & the eighteen examples in \S\ref{sec:examples}\\
\code{scripts/collisions/} & twelve documented collision scripts\\
\code{videos/} & recorded browser videos; \code{videos/scenes/} the scripts behind them\\
\code{tools/} & the index builder, the verifiers, \code{record\_video.py}\\
\code{evidence/port-7.8.0/} & the logs behind every claim in \S9.5\\
\end{tabular}
\end{center}

\begin{center}
\begin{tabular}{p{5.2cm}p{9.4cm}}
\textbf{document} & \textbf{for}\\\hline
\textbf{\code{SolveIt.md}} / \code{.pdf} & this one: the complete solution guide\\
\code{grammar.md} / \code{.pdf} & every command, the full EBNF, 18 more worked examples\\
\code{ARCHITECTURE.md} & the pinned cross-module contracts\\
\code{PORT\_7.8.0\_PROVENANCE.md} & what the 7.8.0 upgrade changed, and the evidence it changed nothing else\\
\code{collision\_detection.md} / \code{.pdf} & the collision science, with 12 example scripts\\
\code{scene\_info.md} / \code{.pdf} & the graphics window, and a survey of seven other simulators\\
\code{physical\_object\_simulator.md} / \code{.pdf} & the predecessor user guide, with 14 further examples\\
\code{special\_functions.md} & the special-function library in detail\\
\code{CLAUDE.md} & the working rules for anyone modifying this repository\\
\code{index\_of\_entities.html} & a browsable catalogue of every named entity in the repository\\
\end{tabular}
\end{center}

\vspace{1em}
\begin{center}
\emph{Everything above was produced by the program in this repository.
Every command shown can be typed. Every number shown can be
reproduced.}
\end{center}

\end{document}
